\section{Model-development history}
\label{app:mutation-history}

The early-epoch chronogeometric model adds one early-epoch degree of
freedom $A_{E}$ to the existing 22-parameter late-time family and
jointly refits all 23 parameters to the unchanged likelihood; the
23-parameter best fit is subsequently frozen. The 10-mode late-time
basis matrix definition is held fixed while its coefficients are refit.
Internal labels (M24, M24B, and so on) name dated fitting stages and are
defined on first use below; they carry no physical meaning beyond
provenance.

{\footnotesize\setlength{\tabcolsep}{4pt}
\begin{longtable}{>{\raggedright\arraybackslash}p{0.10\textwidth}>{\raggedright\arraybackslash}p{0.27\textwidth}>{\raggedright\arraybackslash}p{0.31\textwidth}>{\raggedright\arraybackslash}p{0.23\textwidth}}
\caption{Model-development history: scientific test, result, and consequence at each stage.}\\
\toprule
Stage & Scientific test & Result & Consequence \\
\midrule
\endfirsthead
\toprule
Stage & Scientific test & Result & Consequence \\
\midrule
\endhead
\midrule
\endfoot
\bottomrule
\endlastfoot
M24 & Whether flexible low-redshift time-current freedom could move a joint BAO plus CMB plus local $H_{0}$ fit. & The joint fit moved, showing the low-redshift current was a live direction. & The next question was whether that solution could survive supernova distances. \\
\midrule
M24B & Whether low-redshift freedom could reconcile BAO, CMB, and the local $H_{0}$ value near $73.04$ in a discovery fit without Pantheon shape information. & The fit reached approximately $73.04$ using only BAO, CMB, and the local $H_{0}$ constraint. This stage did not test supernova survival. & The next question was supernova survival, which required the full supernova shape. \\
\midrule
M24C / M24E & Whether the high-$H_{0}$ solution survived the full Pantheon+ distance shape, and whether observer-frame and ladder accounting were correct. & The full supernova shape pulled the high-$H_{0}$ solution toward approximately $69$. Frame and ladder checks fixed bookkeeping but did not restore $73$. & The analysis moved to an exact joint likelihood with explicit spatial geometry. \\
\midrule
M25 & Which epoch and geometry family to carry forward, using the exact joint likelihood with the spatial $\Xi$ mapping. & The epoch family and spatial mapping were constructed and run end to end. & Resolution of the epoch basis remained to be tested. \\
\midrule
M26 & Whether variable-epoch resolution improved the exact joint fit. & Variable epochs improved the joint fit but the Hubble constant stayed near $69.2$. & Late-time temporal freedom was then tested to saturation. \\
\midrule
G10 / M26A & Whether additional late-time temporal freedom could reach $73$. & Modes G4 through G10 saturated: the joint fit stopped improving and $H_{0}$ did not reach $73$. Late temporal freedom was therefore not the missing lever. & Attention shifted to reachability diagnostics and then to an early-epoch lever. \\
\midrule
R131 & Whether high-$H_{0}$ basins were reachable by the optimizer, using a penalty homotopy toward $73.04$ rescored on the real likelihood. & The homotopy was an optimizer-only diagnostic that mapped reachable basins when rescored on the unchanged DESI plus non-ladder likelihood. It was never part of the M28 model. & Its end points were used only as warm starts for subsequent physical fits. \\
\midrule
M27 & Whether differential photon-versus-matter coupling helped, via a one-parameter ratio $\eta$ near $0.988$. & The gain over the undifferentiated fit was too small to matter, with $\eta$ remaining within about one percent of unity. & The differential-coupling direction was closed. \\
\midrule
Ceiling & Whether rescaling the acoustic-ruler amplitude removed the conflict. & A reduced ruler of approximately $7.7$\% removed the conflict and set an upper bound on what ruler physics could explain. & The task became finding a physical early-epoch field that reproduced that amplitude. \\
\midrule
M28 & Whether a physical early-epoch amplitude $A_{E}$, added to the model with all 23 parameters jointly refit to the unchanged likelihood, reproduced the ruler reduction. & The jointly refit 23-parameter fit reached $H_{0}$ near $73$ with $\Delta\chi^{2}\sim-28.6$ relative to the late-time baseline. & The fit required a forward-model and ruler-accounting audit before freezing. \\
\midrule
M28R1 & Forward-model and ruler-accounting audit of the M28 fit on the same likelihood and domain, including rebased calibration and corrected ruler bookkeeping. & The audit confirmed an interior calibrated solution with consistent ruler accounting and unchanged total fit quality. & The audited 23-parameter fit was frozen as the reported model. \\
\end{longtable}
}

\subsection{M24 and M24B: low-redshift discovery without supernova shape}
\label{app:history-m24-m24b}

M24 asked whether flexible low-redshift time-current freedom could move
a joint fit to baryon acoustic oscillations, the cosmic microwave
background, and the local $H_{0}$ constraint. The joint fit moved, which
established the low-redshift current as a direction worth pursuing.

M24B then asked how far that freedom could go in a discovery
configuration without Pantheon shape information. Using only BAO, the
CMB prior, and the local $H_{0}$ value near $73.04$, the fit reconciled
those three inputs near $73.04$. That configuration contained no
supernova distance-shape term, so it established reachability under
restricted data rather than survival under the full data. The open
question was therefore explicit: whether the solution survived the full
supernova Hubble diagram.

\subsection{M24C and M24E: full supernova shape and observer-frame accounting}
\label{app:history-m24c-m24e}

M24C and M24E imposed the full Pantheon+ distance shape together with
the SH0ES ladder rows. The high-$H_{0}$ solution did not survive: the
supernova shape pulled the fit back toward approximately $69$. This was
a data result, not a bookkeeping error.

The same stages rebuilt the observer-frame and ladder machinery to make
sure the pull was not an artifact of frame conventions. M24E
established the homogeneous-lapse observer frame with no free endpoint
factor: photon arrival uses the same lapse as emission with no
additional free rescaling at the observer. M25 then added the spatial
$\Xi$ mapping
\begin{equation}
1+z_{\mathrm{obs}} = (1+z_{\mathrm{bg}})\,\Xi_{0}/\Xi(z),
\end{equation}
so observed redshift folds the spatial field while the background
expansion path uses $z_{\mathrm{bg}}$. DESI uses
$z_{\mathrm{obs}}\to z_{\mathrm{bg}}\to D_{M},D_{H},D_{V}\to D/r_{d}$;
Pantheon+SH0ES uses $z_{\mathrm{HD}}$ for $D_{M}$ with the
$z_{\mathrm{HD}}$/$z_{\mathrm{HEL}}$ prefactor
$D_{L}=(1+z_{\mathrm{HEL}})D_{M}(z_{\mathrm{HD}})$.
Correcting the frame and ladder accounting left the approximately $69$
result intact, so the program moved to an exact joint likelihood with
explicit geometry.

\subsection{M25: exact joint likelihood and spatial geometry}
\label{app:history-m25}

M25 rebuilt the analysis around the exact joint likelihood with the
spatial $\Xi$ mapping above. This replaced approximate or staged
objectives with a single joint objective evaluated on all data, so that
epoch, geometry, and calibration choices competed on the same footing.
The M25 family ran end to end and provided the backbone for resolution
tests. The remaining choice was how finely to resolve the epoch
dependence.

\subsection{M26 and G10/M26A: late-time freedom tested to saturation}
\label{app:history-m26-g10}

M26 varied the epoch resolution of the late-time basis. Finer or
shifted epochs improved the joint fit, but the Hubble constant stayed
near $69.2$ rather than moving toward $73$.

G10/M26A then pushed late-time temporal freedom to saturation by
extending the mode sequence through G4, G6, G8, and G10 and retaining
the first ten modes. The joint fit stopped improving across G8 to G10,
and $H_{0}$ still did not reach $73$. The conclusion was that late-time
temporal freedom, however resolved, was not the lever that reconciled
the local $H_{0}$ value with BAO, CMB, and supernova distances. The next
step had to be an early-epoch lever.

\subsection{R131: optimizer reachability diagnostic}
\label{app:history-r131}

R131 asked a narrower optimizer question: whether high-$H_{0}$ basins
existed at all and could be reached from penalized starts. It applied a
penalty homotopy pulling toward $73.04$ and then rescored every endpoint
on the real, unpenalized likelihood. The diagnostic mapped reachable
basins on DESI plus non-ladder data when rescored honestly. It changed
no physics, entered no model definition, and was never part of M28. Its
only downstream use was to supply diverse warm starts for subsequent
physical fits.

\subsection{M27: differential photon-matter coupling}
\label{app:history-m27}

M27 tested whether photons and matter could propagate under slightly
different couplings, parameterized by a single ratio $\eta$ fitted near
$0.988$. The joint fit improved by an amount too small to justify the
extension, and $\eta$ stayed within about one percent of unity across
starts. The data therefore did not want differential coupling, and the
direction was closed without further follow-up.

\subsection{Ceiling: ruler-amplitude bound}
\label{app:history-ceiling}

The ceiling test asked what ruler amplitude the data wanted if the
amplitude were freed directly. A ruler reduced by approximately $7.7$\%
removed the conflict between the local $H_{0}$ value and the BAO, CMB,
and supernova inputs. This was a bound on mechanism, not a model: it
showed that a smaller early-epoch sound horizon accounts for the
tension, and it set the amplitude a physical early-epoch mechanism had to reproduce.

\subsection{M28: physical early-epoch field with full joint refit}
\label{app:history-m28}

M28 introduced the physical early-epoch amplitude $A_{E}$ and jointly
refit all 23 parameters --- the 22 late-time coefficients plus $A_{E}$
--- to the unchanged likelihood over the full domain. No law was held
fixed during the fit; the late-time basis matrix definition was fixed
but every coefficient was free alongside $A_{E}$. The jointly refit
solution reproduced the ceiling amplitude, reached $H_{0}$ near $73$,
and improved the joint fit by $\Delta\chi^{2}\sim-28.6$ relative to the
late-time baseline. Because the improvement was large and the new
physics touched the ruler, the fit was held for audit rather than
reported immediately.

\subsection{M28R1 audit and freeze}
\label{app:history-m28r1}

M28R1 was the forward-model and ruler-accounting audit of the M28 fit,
run on the same likelihood and domain. It rebased the calibration,
corrected the ruler bookkeeping, and verified that the calibrated
solution sat at an interior point with consistent accounting and
unchanged total fit quality. With the audit passed, the 23-parameter
best fit was frozen as the predecessor publication vector; it is
superseded by Model~A for the branch comparison
(Sec.~\ref{sec:predecessor}).

\subsection{R7A refit, R7B repair, and the final Model~A}
\label{app:history-r7ab}

The M28R1 vector above is predecessor evidence, not the model compared
against the CMB branch in the main text. Two subsequent stages
produced canonical Model~A
(\path{runs/cmb_reconstruction/frozen/MODEL_A_FINAL.json}). First, the
R7A Boltzmann-closed nested refit (3 outer coordinates plus 20 inner,
$x_{1}$ and $x_{2}$ still free) moved the background total only
slightly from the predecessor $\MXXVIIITotal$. Second, the R7B repair
scaled the late spatial
$\Xi$ mapping to zero (L04 center, $x_{1}=x_{2}=0$) without refit,
identifying that mapping as the CMB poison while preserving the
low-redshift fit at DESI+SN $\chi^{2}=\MXXVIIIBTwoDesiSnA$. Final Model~A therefore carries
$H_{0,\mathrm{phys}}=\MXXVIIIBTwoHZeroA$\kmsMpc{} at background
$\MXXVIIIBTwoHZeroBgA$\kmsMpc{},
$\Omega_{m}=\MXXVIIIBTwoOmegaA$, $A_{E}=\MXXVIIIBTwoAE$, and
$r_{\mathrm{drag}}=\MXXVIIIBTwoRdA$~Mpc
(Table~\ref{tab:model-a}): 21 effectively fitted coordinates with two
fixed at zero, distinct from the 23-coordinate predecessor above.

\paragraph{Search-envelope history.}
The M26 canonical $h_{\mathrm{rd}}$ box $[90,115]$~Mpc was a heuristic
optimizer envelope (PROJECT Bounds vs Physical gates), rebased to the
final M28 domain $[70,120]$~Mpc during the M28R1 audit. The bound
history lives here so that Appendix~C states the final domain only.
